% Generated draft: requires mathematical and source review.
\documentclass[11pt]{article}
\usepackage[margin=25mm]{geometry}
\usepackage{fontspec}
\setmainfont{Latin Modern Roman}
\usepackage{amsmath,amssymb,mathtools,unicode-math}
\setmathfont{Latin Modern Math}
\usepackage{enumitem,needspace,xurl,hyperref}
\setlength{\parindent}{0pt}
\setlength{\parskip}{5pt}
\setlength{\emergencystretch}{3em}
\newcommand{\N}{\mathbb{N}}
\newcommand{\Z}{\mathbb{Z}}
\newcommand{\Q}{\mathbb{Q}}
\newcommand{\R}{\mathbb{R}}
\newcommand{\C}{\mathbb{C}}
\newcommand{\F}{\mathbb{F}}
\newcommand{\A}{\mathbb{A}}
\newcommand{\PP}{\mathbb{P}}
\newcommand{\E}{\mathbb{E}}
\newcommand{\eps}{\varepsilon}
\newcommand{\dd}{\,\mathrm{d}}
\begin{document}
\section*{UnsolvedMath ID: 30002527}
\textbf{Source problem code:} catalogue entry
\subsection{Definitions and mathematical statement}
Let \(\mathbb H^3\) be hyperbolic three-space with distance \(d\), and let \(\Gamma<\operatorname{PO}(3,1)\) be a discrete torsion-free subgroup acting by hyperbolic isometries. For \(x\in\mathbb H^3\) and \(\gamma\in\Gamma\), define the closed Dirichlet cell
\[
P_\gamma(x)=\{z\in\mathbb H^3:d(z,\gamma x)\le d(z,\eta x)\text{ for every }\eta\in\Gamma\}.
\]
The Dirichlet tiling \(D_x\) consists of these cells and their faces. For a face \(F\) of this tiling, define its number of incident three-dimensional cells by
\[
n_x(F)=\#\{\gamma\in\Gamma:F\subseteq P_\gamma(x)\}.
\]
The supplied question asks whether, for every such \(\Gamma\), the tiling \(D_x\) is simple for a generic base point \(x\in\mathbb H^3\). If the source's word ``facets'' denotes the maximal three-dimensional cells of the tiling, its stated incidence requirements become
\[
\begin{aligned}
n_x(e)&=3 &&\text{for every one-dimensional face }e\text{ of }D_x,\\
n_x(v)&=4 &&\text{for every zero-dimensional face }v\text{ of }D_x.
\end{aligned}
\]
Here the faces under consideration lie in \(\mathbb H^3\); no compactification or conditions at ideal vertices are supplied.

Two qualifications remain unresolved. First, the source does not specify the meaning of ``generic,'' so it does not determine whether the exceptional base points must have measure zero, form a meagre set, or satisfy another condition. Second, ``facet'' can instead mean a two-dimensional face of a three-dimensional cell. Counting those faces is different from counting incident three-dimensional cells. The displayed formulation is therefore conditional on the maximal-cell interpretation; the supplied wording alone does not select a fully precise formulation.
\subsection{Short English statement}
For every discrete torsion-free group of hyperbolic three-space isometries, ask whether a generic base point produces a Dirichlet tiling with three tiles incident to every edge and four tiles incident to every vertex. This interpretation requires clarification of the source's word ``facets'' and of ``generic.'' The supplied record labels the problem open; this draft does not independently verify that status.
\subsection{Sources}
\begin{itemize}
\item Supplied local UnsolvedMath registry record, permanent ID \(30002527\), title ``Generic Simplicity of Hyperbolic Dirichlet Tilings,'' source problem code \texttt{catalogue entry}, in the collection ``Oberwolfach Reports Open Problems.''
\item Supplied source citation: \emph{Mini-Workshop: K\"ahler Groups} (2014), Oberwolfach Reports, pp.~531--553, DOI \texttt{10.4171/owr/2014/09}. \url{https://doi.org/10.4171/owr/2014/09}. The record attributes the problem to this workshop report and supplies no individual proposer.
\item Supplied literature reference: ``Kapovich, 2013.'' \url{https://doi.org/10.1007/s00222-013-0453-4}. No article title or specific theorem supporting the assessment is supplied.
\item Supplied literature reference: ``Oberwolfach Report 9, 2014.'' \url{https://doi.org/10.4171/owr/2014/09}.
\item Literature assessment dated exactly \texttt{2026-08-21}: the supplied curation labels the unrestricted generic simplicity conjecture open and reports only partial results under extra hypotheses. This is supplied dated curation, not an independently verified literature assessment; the record explicitly requests independent verification of time-sensitive status and cited claims.
\item The supplied record labels its statement \texttt{corrected\_verified}, explaining that mathematical content was preserved while notation, typography, or omitted definitions were repaired. This is a registry qualification, not mathematical verification of the present draft.
\end{itemize}
\end{document}
